\section{Proof of equivalence}
We prove
\begin{equation}
    E\left\{\lim_{T\to\infty}\frac{1}{2T}\int_{-T}^{T} z(t)^Tz(t)\,dt\right\}
    \;=\;
    \operatorname{tr}E\{z(t)z(t)^T\}
\end{equation}

\begin{equation}
    \implies
    \frac{1}{2\pi}\int_{-\infty}^{\infty}\operatorname{tr}\big[F(j\omega)F(j\omega)^H\big]\,d\omega
    \;=\;
    \|F\|_2^2,
\end{equation}

where $z(t) = Fl(P,K)w(t) \triangleq F(s)w(t)$, and $w(t)$ is white noise of unit intensity,

\[
    E\{w(t)w(\tau)^T\} = I\,\delta(t-\tau)
\]

\bigskip
\noindent The proof proceeds in four steps.

\section*{Step 1: Scalar quadratic form as a trace}

For any real vector $v \in \mathbb{R}^n$,
\[
    v^Tv = \sum_i v_i^2 = \operatorname{tr}(vv^T),
\]
since $vv^T$ is the matrix with $(vv^T)_{ii} = v_i^2$, and the trace sums the diagonal entries. Applying this with $v = z(t)$:
\[
    z(t)^Tz(t) = \operatorname{tr}\big(z(t)z(t)^T\big).
\]

\section*{Step 2: Trace commutes with time-average and expectation}

The trace operator, the time-integral, and the expectation $E\{\cdot\}$ are all linear operators, hence they commute with one another. Substituting Step 1 into the left-hand side of (9.31):
\[
    E\left\{\lim_{T\to\infty}\frac{1}{2T}\int_{-T}^{T} z(t)^Tz(t)\,dt\right\}
    = E\left\{\lim_{T\to\infty}\frac{1}{2T}\int_{-T}^{T} \operatorname{tr}\big(z(t)z(t)^T\big)\,dt\right\}.
\]
Pulling $\operatorname{tr}(\cdot)$ outside the (linear) time-average and expectation:
\[
    = \operatorname{tr}\left(E\left\{\lim_{T\to\infty}\frac{1}{2T}\int_{-T}^{T} z(t)z(t)^T\,dt\right\}\right).
\]

\section*{Step 3: Ergodicity collapses the time average to the ensemble average}

Since $z(t) = F(s)w(t)$ is the output of a stable LTI system driven by stationary white noise, $z(t)$ is a stationary stochastic process. For such processes the ergodic theorem states that the time average of a (sufficiently well-behaved) function of the process equals its ensemble average with probability one:
\[
    \lim_{T\to\infty}\frac{1}{2T}\int_{-T}^{T} z(t)z(t)^T\,dt = E\{z(t)z(t)^T\} \quad \text{(a deterministic matrix, independent of } t\text{)}.
\]
Substituting this in:
\[
    \operatorname{tr}\left(E\Big\{\underbrace{E\{z(t)z(t)^T\}}_{\text{already deterministic}}\Big\}\right) = \operatorname{tr}\, E\{z(t)z(t)^T\},
\]
since applying $E\{\cdot\}$ to an already-deterministic quantity leaves it unchanged. This establishes the first equality in (9.31):
\[
    E\left\{\lim_{T\to\infty}\frac{1}{2T}\int_{-T}^{T} z(t)^Tz(t)\,dt\right\} = \operatorname{tr}\, E\{z(t)z(t)^T\}.
\]

\section*{Step 4: Introducing the system dynamics and Parseval's theorem}

\subsection*{4.1 Express $z$ as a convolution}

Let $f(\theta)$ denote the impulse response matrix of the stable system $F(s)$. Then
\[
    z(t) = \int_0^\infty f(\theta_1)\,w(t-\theta_1)\,d\theta_1.
\]

\subsection*{4.2 Compute $E\{z(t)z(t)^T\}$}

\[
    E\{z(t)z(t)^T\}
    = E\left\{\int_0^\infty f(\theta_1)w(t-\theta_1)\,d\theta_1 \int_0^\infty w(t-\theta_2)^Tf(\theta_2)^T\,d\theta_2\right\}.
\]
Moving $E\{\cdot\}$ inside the (linear) double integral:
\[
    = \int_0^\infty\!\!\int_0^\infty f(\theta_1)\,E\{w(t-\theta_1)w(t-\theta_2)^T\}\,f(\theta_2)^T\,d\theta_1\,d\theta_2.
\]
By the white-noise property (9.30), with $\tau_1 = t-\theta_1,\ \tau_2 = t-\theta_2$:
\[
    E\{w(t-\theta_1)w(t-\theta_2)^T\} = I\,\delta\big((t-\theta_1)-(t-\theta_2)\big) = I\,\delta(\theta_2-\theta_1).
\]
Substituting and using the sifting property of the Dirac delta, $\int_0^\infty g(\theta_2)\,\delta(\theta_2-\theta_1)\,d\theta_2 = g(\theta_1)$, the double integral collapses to a single integral:
\[
    E\{z(t)z(t)^T\} = \int_0^\infty f(\theta_1)f(\theta_1)^T\,d\theta_1 = \int_0^\infty f(\theta)f(\theta)^T\,d\theta.
\]

\subsection*{4.3 Apply the trace}

\[
    \operatorname{tr}\,E\{z(t)z(t)^T\} = \operatorname{tr}\int_0^\infty f(\theta)f(\theta)^T\,d\theta = \int_0^\infty \operatorname{tr}\big(f(\theta)f(\theta)^T\big)\,d\theta,
\]
using linearity of trace to move it inside the integral.

\subsection*{4.4 Identify the $H_2$ norm and apply Parseval's theorem}

By the cyclic property of trace, $\operatorname{tr}(AB) = \operatorname{tr}(BA)$, applied with $A = f(\theta)$ and $B = f(\theta)^T$:
\[
    \operatorname{tr}\big(f(\theta)f(\theta)^T\big) = \operatorname{tr}\big(f(\theta)^Tf(\theta)\big),
\]
so the integral matches the definition of the (squared) $H_2$ norm of the impulse response,
\[
    \int_0^\infty \operatorname{tr}\big(f(\theta)^Tf(\theta)\big)\,d\theta = \|F\|_2^2.
\]
By Parseval's theorem, this time-domain integral equals the corresponding frequency-domain integral:
\[
    \int_0^\infty \operatorname{tr}\big(f(\theta)f(\theta)^T\big)\,d\theta
    = \frac{1}{2\pi}\int_{-\infty}^{\infty} \operatorname{tr}\big[F(j\omega)F(j\omega)^H\big]\,d\omega.
\]

\section*{Conclusion}

Combining Steps 1--4:
\[
    E\left\{\lim_{T\to\infty}\frac{1}{2T}\int_{-T}^{T} z(t)^Tz(t)\,dt\right\}
    = \operatorname{tr}\,E\{z(t)z(t)^T\}
    = \frac{1}{2\pi}\int_{-\infty}^{\infty}\operatorname{tr}\big[F(j\omega)F(j\omega)^H\big]\,d\omega
    = \|F\|_2^2,
\]
which is precisely equation (9.31). $\blacksquare$
